\frfilename{mt245.tex} \versiondate{25.3.06} \copyrightdate{1995}
\def\halfarrow{\rightharpoonup}

\loadeusm

\def\chaptername{Ruang fungsi} \def\sectionname{Konvergensi dalam ukuran}

\newsection{245}

Kini saya membahas suatu topologi penting dan menarik pada ruang
$\eusm L^0$ dan $L^0$. Saya mulai dengan definisinya (245A), lalu meninjau
sifat-sifat yang serupa dengan sifat ruang $L^p$ untuk $p\ge 1$ (245D-245E).
Dalam 245G-245J saya menjelaskan hubungan-hubungan yang paling berguna antara
topologi ini dan topologi norma pada ruang $L^p$. Untuk ruang
$\sigma$-hingga, topologi ini dapat dimetriskan (245Eb), dan konvergensi
barisan dapat dicirikan melalui konvergensi titik demi titik suatu barisan
fungsi (245K-245L).

\leader{245A}{Definisi} Misalkan $(X,\Sigma,\mu)$ adalah ruang ukur.

\header{245Aa}{\bf (a)} Untuk setiap himpunan terukur $F\subseteq X$ yang
berukuran berhingga, terdapat fungsional $\tau_F$ pada
$\eusm L^0=\eusm L^0(\mu)$ yang didefinisikan oleh

\Centerline{$\tau_F(f)=\int|f|\wedge\chi F$}

\noindent untuk setiap $f\in\eusm L^0$. \cmmnt{(Integral ini ada di $\Bbb R$
karena $|f|\wedge\chi F$ termasuk dalam $\eusm L^0$ dan didominasi oleh
fungsi terintegralkan $\chi F$).} Selain itu,
$\tau_F(f+g)\le\tau_F(f)+\tau_F(g)$ setiap kali $f$, $g\in\eusm L^0$.
\prooflet{\Prf\ Cukup perhatikan bahwa

\Centerline{$\min(|(f+g)(x)|,\chi F(x))
\le\min(|f(x)|,\chi
F(x))+\min(|g(x)|,\chi F(x))$}

\noindent untuk setiap $x\in\dom f\cap\dom g$, yang berlaku untuk hampir setiap $x\in X$.\
\QeD} Akibatnya, dengan menetapkan $\rho_F(f,g)=\tau_F(f-g)$,\cmmnt{ kita
mempunyai

\Centerline{$\rho_F(f,h)=\tau_F((f-g)+(g-h))
\le\tau_F(f-g)+\tau_F(g-h)=\rho_F(f,g)+\rho_F(g,h)$,}

\Centerline{$\rho_F(f,g)=\tau_F(f-g)\ge 0$,}

\Centerline{$\rho_F(f,g)=\tau_F(f-g)=\tau_F(g-f)=\rho_F(g,f)$}

\noindent untuk semua $f$, $g$, $h\in\eusm L^0$; dengan kata lain, } $\rho_F$ adalah
pseudometrik pada $\eusm L^0$.

\header{245Ab}{\bf (b)} Keluarga

\Centerline{$\{\rho_F:F\in\Sigma,\,\mu F<\infty\}$}

\noindent mendefinisikan suatu topologi pada $\eusm L^0$\cmmnt{ (2A3F)};
saya akan menyebutnya topologi {\bf konvergensi dalam ukuran} pada
$\eusm L^0$.

\header{245Ac}{\bf (c)} Jika $f$, $g\in\eusm L^0$ dan $f\eae g$, maka\cmmnt{
$|f|\wedge\chi F\eae|g|\wedge\chi F$ dan} $\tau_F(f)=\tau_F(g)$, untuk setiap
himpunan $F$ berukuran berhingga. Akibatnya kita mempunyai fungsional
$\bar\tau_F$ pada $L^0=L^0(\mu)$ yang didefinisikan melalui

\Centerline{$\bar\tau_F(f^{\ssbullet})=\tau_F(f)$}

\noindent setiap kali $f\in\eusm L^0$, $F\in\Sigma$ dan $\mu F<\infty$.
Fungsional ini menghasilkan pseudometrik $\bar\rho_F$ yang dapat dinyatakan
dengan salah satu rumus

\Centerline{$\bar\rho_F(u,v)=\bar\tau_F(u-v)$,
\quad$\bar\rho_F(f^{\ssbullet},g^{\ssbullet})=\rho_F(f,g)$}

\noindent untuk $u$, $v\in L^0$, $f$, $g\in\eusm L^0$ dan $F$ dengan ukuran
berhingga. Keluarga pseudometrik ini mendefinisikan topologi
{\bf konvergensi dalam ukuran} pada $L^0$.

\spheader 245Ad Selanjutnya, suatu barisan (dalam $\eusm L^0$
atau $L^0$) disebut {\bf konvergen dalam ukuran} jika barisan itu konvergen dalam
topologi konvergensi dalam ukuran\cmmnt{ (dalam arti 2A3M)}.

\cmmnt{ \leader{245B}{Catatan (a)} Tentu saja topologi pada $\eusm L^0$ dan
$L^0$ berkaitan seerat mungkin. Topologi pada $L^0$ tidak hanya
merupakan hasil bagi topologi pada $\eusm L^0$ (yaitu, himpunan $G\subseteq
L^0$ terbuka jika $\{f:f^{\ssbullet}\in G\}$ terbuka di $\eusm L^0$), tetapi
setiap himpunan terbuka di $\eusm L^0$ juga merupakan prapeta, melalui peta
hasil bagi, dari suatu himpunan terbuka di $L^0$.


\header{245Bb}{\bf (b)} Perhatikan pula bahwa jika $F_0,\ldots,F_n$
adalah himpunan terukur berukuran berhingga dengan gabungan $F$, maka, dalam
notasi 245A, $\tau_{F_i}\le\tau_{F}$ untuk setiap $i$;   ini berarti bahwa
himpunan $G\subseteq \eusm L^0$ terbuka untuk topologi konvergensi dalam
ukuran jika dan hanya jika, untuk setiap $f\in G$, terdapat himpunan $F$
berukuran berhingga dan $\delta>0$ sedemikian sehingga

\Centerline{$\rho_F(g,f)\le\delta\Longrightarrow g\in G$.}

\noindent Demikian pula, himpunan $G\subseteq L^0$ terbuka untuk topologi
konvergensi dalam ukuran jika dan hanya jika, untuk setiap $u\in G$, terdapat
himpunan $F$ berukuran berhingga dan $\delta>0$ sehingga

\Centerline{$\bar\rho_F(v,u)\le\delta\Longrightarrow v\in G$.}

\header{245Bc}{\bf (c)} Istilah `topologi konvergensi dalam ukuran' cukup
sesuai dengan penggunaan standar ketika $(X,\Sigma,\mu)$ berhingga total.
Namun {\bf hati-hati!} Istilah `topologi konvergensi dalam ukuran' juga
digunakan untuk topologi yang ditentukan oleh metrik 245Ye di bawah, bahkan
ketika $\mu X=\infty$. Saya juga pernah menjumpai istilah
{\bf konvergensi lokal dalam ukuran} untuk topologi 245A. Kebanyakan penulis mengabaikan
ruang yang tidak $\sigma$-hingga dalam konteks ini. Namun saya berpendapat bahwa
245D-245E di bawah cukup menarik sehingga perlu diperluas ke kasus tersebut.
}%akhir komentar

\leader{245C}{Konvergensi titik demi titik}\cmmnt{ Topologi konvergensi dalam
ukuran hampir dapat didefinisikan melalui `konvergensi titik demi titik', yang
merupakan salah satu akar teori ukuran. Kesesuaiannya paling erat pada ruang
ukur $\sigma$-hingga (lihat 245K), tetapi dalam kasus umum pun masih ada
hubungan yang sangat penting, sebagai berikut. } Misalkan $(X,\Sigma,\mu)$
adalah ruang ukur, dan tulis $\eusm
L^0=\eusm L^0(\mu)$, $L^0=L^0(\mu)$.

\spheader 245Ca Jika $\sequencen{f_n}$ adalah barisan dalam $\eusm L^0$ yang
konvergen hampir di mana-mana ke $f\in\eusm L^0$, maka
$\sequencen{f_n}\to f$ dalam ukuran. \prooflet{\Prf\ Menurut 2A3Mc, cukup
dibuktikan bahwa $\lim_{n\to\infty}\rho_F(f_n,f)=0$ setiap kali $\mu
F<\infty$. Tetapi $\sequencen{|f_n-f|\wedge\chi F}$ konvergen ke $0$ hampir
di mana-mana dan didominasi oleh fungsi terintegralkan $\chi F$, sehingga,
menurut Teorema Konvergensi Terdominasi Lebesgue,

\Centerline{$\lim_{n\to\infty}\rho_F(f_n,f)
=\lim_{n\to\infty}\int|f_n-f|\wedge\chi F=0$.   \Qed}}

\spheader 245Cb \cmmnt{ Untuk merumuskan hasil terkait yang dapat diterapkan
pada $L^0$, kita memerlukan konsep berikut.   Jika $\sequencen{f_n}$,
$\sequencen{g_n}$ adalah barisan dalam $\eusm L^0$ sehingga
$f_n^{\ssbullet}=g_n^{\ssbullet}$ untuk setiap $n$, sedangkan $f$, $g\in\eusm L^0$
memenuhi $f^{\ssbullet}=g^{\ssbullet}$, dan
$\sequencen{f_n}\to f$ hampir di mana-mana, maka
$\sequencen{g_n}\to g$ hampir di mana-mana\prooflet{, karena

$$\eqalign{\{x:x\in\dom f&\cap\dom g\cap\bigcap_{n\in\Bbb N}\dom
f_n\cap\bigcap_{n\in\Bbb N}g_n,\cr
&g(x)=f(x)=\lim_{n\to\infty}f_n(x),\,f_n(x)
=g_n(x)\Forall n\in \Bbb N\}\cr}$$

\noindent koterabaikan}. Akibatnya diperoleh definisi bagi
barisan dalam $L^0$; untuk} \dvro{Untuk}{}
$f$, $f_n\in\eusm L^0$, barisan $\sequencen{f_n^{\ssbullet}}$ dikatakan {\bf
konvergen secara ordo*}, atau {\bf ordo*-konvergen}, ke
$f^{\ssbullet}$ jika dan hanya jika $f\eae\lim_{n\to\infty}f_n$.
Dalam\cmmnt{ hal ini, tentu saja, $\sequencen{f_n}\to f$ dalam ukuran.
Jadi, dalam} $L^0$, barisan $\sequencen{u_n}$ yang ordo*-konvergen ke
$u\in L^0$ juga konvergen dalam ukuran ke $u$.

\cmmnt{\medskip

\noindent{\bf Catatan} Uraian alternatif tentang konvergensi
ordo diberikan dalam 245Xc; kondisi (iii)-(vi) di sana disajikan dalam bentuk yang sesuai
dengan struktur yang lebih umum. }%akhir komentar

\spheader 245Cc\cmmnt{ Untuk contoh tipikal barisan yang konvergen dalam
ukuran tanpa konvergen secara ordo, pertimbangkan contoh berikut. } Ambil
$\mu$ sebagai ukuran Lebesgue pada $[0,1]$, dan tetapkan $f_n(x)=2^m$ jika
$x\in[2^{-m}k,2^{-m}(k+1)]$, $0$ sebaliknya, di mana $k=k(n)\in\Bbb N$,
$m=m(n)\in\Bbb N$ ditentukan oleh syarat $n+1=2^{m}+k$ dan
$0\le k<2^m$. Maka $\sequencen{f_n}\to 0$ untuk topologi konvergensi
dalam ukuran\cmmnt{ (karena $\rho_F(f_n,0)\le
2^{-m}$ jika $F\subseteq[0,1]$
terukur dan $2^m-1\le n$)}, meskipun $\sequencen{f_n}$ tidak konvergen ke
$0$ hampir di mana-mana\cmmnt{ (memang, $\limsup_{n\to\infty}f_n=\infty$ di
mana-mana)}.

\leader{245D}{Proposisi} Misalkan $(X,\Sigma,\mu)$ adalah sebarang ruang ukur.

(a) Topologi konvergensi dalam ukuran merupakan topologi ruang linier pada
$L^0=L^0(\mu)$.

(b) Peta $\vee$, $\wedge:L^0\times L^0\to L^0$, dan $u\mapsto|u|$, $u\mapsto
u^+$, $u\mapsto u^-:L^0\to L^0$ semuanya kontinu.

(c) Peta $\times:L^0\times L^0\to L^0$ kontinu.

(d) Untuk setiap fungsi kontinu $h:\Bbb R\to\Bbb R$, fungsi terkait $\bar
h:L^0\to L^0$\cmmnt{ (241I)} adalah kontinu.

\proof{{\bf (a)} Kuncinya ialah bahwa fungsional $\bar\tau_F$,
sebagaimana didefinisikan dalam 245Ac, merupakan F-seminorma dalam arti 2A5B.
\Prf\ Ambil himpunan $F\in\Sigma$ berukuran berhingga. Dalam 245Aa telah dicatat bahwa

\Centerline{$\tau_F(f+g)\le\tau_F(f)+\tau_F(g)$ untuk semua $f$, $g\in\eusm
L^0$,}

\noindent jadi

\Centerline{$\bar\tau_F(u+v)\le\bar\tau_F(u)+\bar\tau_F(v)$ untuk semua $u$,
$v\in L^0$.}

\noindent Selanjutnya,

$$\eqalignno{\bar\tau_F(cu)&\le\bar\tau_F(u)\text{ setiap kali }u\in L^0
  \text{
dan }|c|\le 1&\text{(*)}\cr}$$

\noindent karena $|cf|\wedge\chi F\leae|f|\wedge\chi F$ setiap kali $f\in\eusm
L^0$ dan $|c|\le 1$. Terakhir, jika diberikan $u\in L^0$ dan $\epsilon>0$,
ambil $f\in\eusm L^0$ sedemikian sehingga $f^{\ssbullet}=u$. Maka

\Centerline{$\lim_{n\to\infty}|2^{-n}f|\wedge\chi F=0$ hampir di mana-mana,}

\noindent sehingga, menurut Teorema Konvergensi Terdominasi Lebesgue,

\Centerline{$\lim_{n\to\infty}\bar\tau_F(2^{-n}u)
=\lim_{n\to\infty}\int|2^{-n}f|\wedge\chi F=0$,}

\noindent dan terdapat $n$ sehingga $\bar\tau_F(2^{-n}u)\le\epsilon$.
Dari (*) di atas diperoleh bahwa $\bar\tau_F(cu)\le\epsilon$ setiap
kali $|c|\le 2^{-n}$. Karena $\epsilon$ sembarang, $\lim_{c\to
0}\bar\tau_F(cu)=0$ untuk setiap $u\in L^0$;  yang merupakan kondisi ketiga di
2A5B.\ \Qed

Menurut 2A5B, topologi yang ditentukan oleh
$\bar\tau_F$ merupakan topologi ruang linier.

\medskip

{\bf (b)} Untuk $u$, $v\in L^0$, $||u|-|v||\le|u-v|$, jadi
$\bar\rho_F(|u|,|v|)\le\bar\rho_F(u,v)$ untuk setiap himpunan $F$ berukuran
berhingga. Menurut 2A3H, $|\,\,|:L^0\to L^0$ kontinu. Selanjutnya,

\Centerline{$u\vee v=\Bover12(u+v+|u-v|)$, \quad$u\wedge
v=\Bover12(u+v-|u-v|)$,}

\Centerline{$u^+=u\vee 0$, \quad$u^-=(-u)\vee 0$.}

\noindent Karena penjumlahan dan pengurangan adalah kontinu, maka $\vee$,
$\wedge$, $^+$ dan $^-$ adalah kontinu.

\medskip

{\bf (c)} Ambil $u_0$, $v_0\in L^0$, $F\in\Sigma$ suatu himpunan berukuran
berhingga, dan $\epsilon>0$. Nyatakan $u_0$ dan $v_0$ masing-masing sebagai
$f_0^{\ssbullet}$ dan $g_0^{\ssbullet}$, dengan $f_0$, $g_0:X\to\Bbb R$
terukur terhadap $\Sigma$ (241Bk). Jika kita menetapkan

\Centerline{$F_m=\{x:x\in F,\,|f_0(x)|+|g_0(x)|\le m\}$,}

\noindent maka $\sequence{m}{F_m}$ merupakan barisan himpunan tak menurun
dengan gabungan $F$; karena itu terdapat $m\in\Bbb N$ dengan $\mu(F\setminus
F_m)\le\bover12\epsilon$. Pilih $\delta>0$ sedemikian sehingga
$(2m+\mu F)\delta^2+2\delta\le\bover12\epsilon$.

Sekarang andaikan $u$, $v\in L^0$ memenuhi
$\bar\rho_F(u,u_0)\le\delta^2$ dan $\bar\rho_F(v,v_0)\le\delta^2$. Pilih
$f$, $g:X\to\Bbb R$ yang terukur dan memenuhi
$f^{\ssbullet}=u$ dan $g^{\ssbullet}=v$. Maka

\Centerline{$\mu\{x:x\in F,\,|f(x)-f_0(x)|\ge\delta\}\le\delta$,
\quad$\mu\{x:x\in F,\,|g(x)-g_0(x)|\ge\delta\}\le\delta$,}

\noindent sehingga

\Centerline{$\mu\{x:x\in F,\,|f(x)-f_0(x)||g(x)-g_0(x)|\ge\delta^2\}
\le
2\delta$}

\noindent dan

\Centerline{$\int_F\min(1,|f-f_0|\times|g-g_0|)
\le\delta^2\mu F+2\delta$.}

\noindent Selain itu,

\Centerline{$|f\times g-f_0\times g_0|
\le|f-f_0|\times|g-g_0| +
|f_0|\times|g-g_0| + |f-f_0|\times|g_0|$,}

\noindent sehingga

$$\eqalignno{\bar\rho_F(u\times v,u_0\times v_0)
&=\int_F\min(1,|f\times
g-f_0\times g_0|)\cr
&\le\Bover12\epsilon+\int_{F_m}\min(1,|f\times
g-f_0\times g_0|)\cr
&\le\Bover12\epsilon+\int_{F_m}\min(1,|f-f_0|\times|g-g_0|
+m|g-g_0|+m|f-f_0|)\cr
&\le\Bover12\epsilon+\int_F\min(1,|f-f_0|\times|g-g_0|)\cr
&\qquad\qquad\qquad\qquad+m\int_F\min(1,|g-g_0|)
+m\int_F\min(1,|f-f_0|)\cr
&\le\Bover12\epsilon+\delta^2\mu
F+2\delta+2m\delta^2
\le\epsilon.\cr}$$

\noindent Karena $F$ dan $\epsilon$ sembarang, $\times$ kontinu di
$(u_0,v_0)$; karena $u_0$ dan $v_0$ sembarang, $\times$ kontinu.

\medskip

{\bf (d)} Ambil $u\in L^0$, $F\in\Sigma$ berukuran berhingga, dan
$\epsilon>0$. Maka terdapat $\delta>0$ sehingga
$\bar\rho_F(\bar h(v),\bar h(u))\le\epsilon$ setiap
kali $\bar\rho_F(v,u)\le\delta$. \Prf\Quer\ Jika tidak, kita dapat menemukan,
untuk setiap $n\in\Bbb N$, suatu $v_n$ sehingga $\bar\rho_F(v_n,u)\le 4^{-n}$
tetapi $\bar\rho_F(\bar h(v_n),\bar h(u))>\epsilon$. Nyatakan $u$ sebagai
$f^{\ssbullet}$ dan $v_n$ sebagai $g_n^{\ssbullet}$ dengan $f$, $g_n:X\to\Bbb
R$ terukur. Tetapkan

\Centerline{$E_n=\{x:x\in F,\,|g_n(x)-f(x)|\ge 2^{-n}\}$}

\noindent untuk setiap $n$. Maka $\bar\rho_F(v_n,u)\ge 2^{-n}\mu E_n$,
jadi $\mu E_n\le 2^{-n}$ untuk setiap $n$, dan $E=\bigcap_{n\in\Bbb
N}\bigcup_{m\ge n}E_m$ terabaikan. Tetapi $\lim_{n\to\infty}g_n(x)=f(x)$
untuk setiap $x\in F\setminus E$, jadi (karena $h$ kontinu)
$\lim_{n\to\infty}h(g_n(x))=h(f(x))$ untuk setiap $x\in F\setminus E$.
Akibatnya (seperti biasa, menurut Teorema Konvergensi Terdominasi Lebesgue)

\Centerline{$\lim_{n\to\infty}\bar\rho_F(\bar h(v_n),\bar h(u))
=\lim_{n\to\infty}\int_F\min(1,|h(g_n(x))-h(f(x))|)\mu(dx)=0$,}

\noindent yang tidak mungkin.\ \Bang\Qed

Menurut 2A3H, $\bar h$ kontinu. }%akhir pembuktian 245D

\cmmnt{\medskip

\noindent{\bf Catatan} Saya tidak dapat mengatakan bahwa topologi
konvergensi dalam ukuran pada $\eusm L^0$ merupakan topologi ruang linier
semata-mata karena (pada definisi yang saya pilih) $\eusm L^0$ pada umumnya
bukan ruang linier.}



\leader{245E}{}\cmmnt{ Kini saya beralih ke teorema utama yang
menghubungkan sifat-sifat ruang linier topologis $L^0(\mu)$ dengan
klasifikasi ruang ukur dalam Bab 21.

\medskip

\noindent}{\bf Teorema} Misalkan $(X,\Sigma,\mu)$ adalah ruang ukur. Misalkan
$\frak T$ adalah topologi konvergensi dalam ukuran pada
$L^0=L^0(\mu)$\cmmnt{, seperti yang dijelaskan dalam 245A}.

(a) $(X,\Sigma,\mu)$ semihingga jika dan hanya jika $\frak T$ Hausdorff.

(b) $(X,\Sigma,\mu)$ $\sigma$-hingga jika dan hanya jika $\frak T$ dapat
dimetriskan.

(c) $(X,\Sigma,\mu)$ dapat dilokalkan jika dan hanya jika $\frak T$ Hausdorff
dan $L^0$ lengkap terhadap $\frak T$.

\wheader{245E}{0}{0}{0}{48pt}

\proof{ Gunakan pseudometrik $\rho_F$ pada
$\eusm L^0=\eusm L^0(\mu)$ dan $\bar\rho_F$ pada $L^0$ yang diuraikan dalam
245A.

\medskip

{\bf (a)(i)} Misalkan $(X,\Sigma,\mu)$ semihingga dan $u$, $v$ dua anggota
berbeda dari $L^0$. Nyatakan keduanya sebagai $f^{\ssbullet}$ dan
$g^{\ssbullet}$, dengan $f$ dan $g$ fungsi terukur dari $X$ ke $\Bbb R$.
Maka $\mu\{x:f(x)\ne g(x)\}>0$; karena $(X,\Sigma,\mu)$ semihingga, terdapat
himpunan $F\in\Sigma$ yang berukuran berhingga
sehingga $\mu\{x:x\in F,\,f(x)\ne g(x)\}>0$. Dengan demikian,

\Centerline{$\bar\rho_F(u,v)=\int_F\min(|f(x)-g(x)|,1)dx>0$}

\noindent (lihat 122Rc). Karena $u$ dan $v$ sembarang, $\frak T$
adalah Hausdorff (2A3L).

\medskip

\quad{\bf (ii)} Misalkan $\frak T$ Hausdorff dan $E\in\Sigma$, $\mu E>0$.
Maka $u=\chi E^{\ssbullet}\ne 0$, sehingga terdapat $F\in\Sigma$ dengan
$\mu F<\infty$ dan $\bar\rho_F(u,0)\ne 0$, yakni $\mu(E\cap F)>0$. Dengan
demikian $E\cap F$ adalah himpunan berukuran berhingga yang tidak terabaikan
dan termuat dalam $E$. Karena $E$ sembarang, $(X,\Sigma,\mu)$ semihingga.

\medskip

{\bf (b)(i)} Misalkan $(X,\Sigma,\mu)$ $\sigma$-hingga. Misalkan
$\sequencen{E_n}$ barisan tak menurun dari himpunan-himpunan berukuran
berhingga yang gabungannya adalah $X$. Tetapkan

$$\bar\rho(u,v)
=\sum_{n=0}^{\infty}\bover{\bar\rho_{E_n}(u,v)}{1+2^n\mu
E_n}$$

\noindent untuk $u$, $v\in L^0$.   Maka $\bar\rho$ adalah metrik pada $L^0$.
\Prf\ Karena setiap $\bar\rho_{E_n}$ merupakan pseudometrik, demikian pula
$\bar\rho$. Jika $\bar\rho(u,v)=0$, nyatakan $u$ sebagai $f^{\ssbullet}$ dan
$v$ sebagai $g^{\ssbullet}$, dengan $f$, $g\in\eusm L^0(\mu)$; maka

\Centerline{$\int|f-g|\wedge\chi E_n=\bar\rho_{E_n}(u,v)=0$,}

\noindent jadi $f=g$ hampir di mana-mana di $E_n$, untuk setiap $n$. Karena
$X=\bigcup_{n\in\Bbb N}E_n$, $f\eae g$ dan $u=v$.\ \Qed\

Jika $F\in\Sigma$, $\mu F<\infty$, dan $\epsilon>0$, ambil $n$ sehingga
$\mu(F\setminus E_n)\le\bover12\epsilon$.   Jika $u$, $v\in L^0$ dan
$\bar\rho(u,v)\le\epsilon/2(1+2^n\mu E_n)$, maka $\bar\rho_F(u,v)\le\epsilon$.
\Prf\ Nyatakan $u$ sebagai $f^{\ssbullet}$ dan $v=g^{\ssbullet}$,
dengan $f$, $g\in\eusm L^0$. Maka

\Centerline{$\int|u-v|\wedge \chi E_n
=\bar\rho_{E_n}(u,v)\le(1+2^n\mu
E_n)\bar\rho(u,v)
\le\Bover{\epsilon}2$,}

\noindent sementara

\Centerline{$\int|f-g|\wedge \chi(F\setminus E_n)
\le\mu(F\setminus
E_n)\le\Bover{\epsilon}2$,}


\noindent jadi

\Centerline{$\bar\rho_F(u,v)=\int|f-g|\wedge\chi F
\le\int|f-g|\wedge\chi
E_n+\int|f-g|\wedge\chi(F\setminus E_n)
\le\Bover{\epsilon}2+\Bover{\epsilon}2
=\epsilon$.  \Qed}

Sebaliknya, jika $\epsilon>0$, ambil $n\in\Bbb N$ sehingga
$2^{-n}\le\bover12\epsilon$;  lalu $\bar\rho(u,v)\le\epsilon$ setiap kali
$\bar\rho_{E_n}(u,v)\le\epsilon/2(n+1)$.

Hal ini menunjukkan bahwa $\bar\rho$ mendefinisikan topologi yang sama dengan
$\bar\rho_F$ (2A3Ib), sehingga $\frak T$, topologi yang ditentukan oleh
$\bar\rho_F$, dapat dimetriskan.

\medskip

\quad{\bf (ii)} Misalkan $\frak T$ dapat dimetriskan. Misalkan $\bar\rho$
suatu metrik yang menentukan $\frak T$. Untuk setiap $n\in\Bbb N$ harus
terdapat himpunan terukur $F_n$ berukuran berhingga dan $\delta_n>0$ sedemikian
rupa sehingga

\Centerline{$\bar\rho_{F_n}(u,0)\le \delta_n\Longrightarrow\bar\rho(u,0)
\le
2^{-n}$.}

\noindent Tetapkan $E=X\setminus\bigcup_{n\in\Bbb N}F_n$. \Quer\ Andaikan $E$
tidak terabaikan, maka $u=\chi E^{\ssbullet}\ne 0$; karena $\bar\rho$
adalah metrik, terdapat $n\in\Bbb N$ sehingga $\bar\rho(u,0)>2^{-n}$. Tetapi


\Centerline{$\mu(E\cap F_n)=\bar\rho_{F_n}(u,0)>\delta_n$.}

\noindent ini bertentangan dengan $E\cap F_n=\emptyset$.\ \Bang

Jadi $\mu E=0<\infty$. Karena $X=E\cup\bigcup_{n\in\Bbb N}F_n$ merupakan
gabungan terhitung dari himpunan-himpunan berukuran berhingga, $\mu$ pun
$\sigma$-hingga.

\medskip

{\bf (c)} Menurut (a), masing-masing hipotesis menjamin bahwa $\mu$ semihingga
dan $\frak T$ Hausdorff.

\medskip

\quad{\bf (i)} Misalkan $(X,\Sigma,\mu)$ dapat dilokalkan, dan misalkan $\Cal F$
suatu filter Cauchy pada $L^0$ (2A5F). Untuk setiap himpunan terukur $F$
berukuran berhingga, pilih barisan $\sequencen{A_n(F)}$ dalam
$\Cal F$ sehingga $\bar\rho_F(u,v)\le 4^{-n}$ untuk setiap $u$, $v\in A_n(F)$,
dan setiap $n$ (2A5G). Pilih $u_{Fn}\in\bigcap_{k\le n}A_k(F)$ untuk setiap
$n$;  lalu $\bar\rho_F(u_{F,n+1},u_{Fn})\le 2^{-n}$ untuk setiap $n$.
Nyatakan setiap $u_{Fn}$ sebagai $f_{Fn}^{\ssbullet}$, dengan $f_{Fn}$
fungsi terukur dari $X$ ke $\Bbb R$. Maka

\Centerline{$\mu\{x:x\in F,\,|f_{F,n+1}(x)-f_{Fn}(x)|\ge 2^{-n}\}
\le
2^n\bar\rho_F(u_{F,n+1},u_{Fn})\le 2^{-n}$}

\noindent untuk setiap $n$. Oleh karena itu, $\sequencen{f_{Fn}}$ harus
konvergen hampir di mana-mana pada $F$. \Prf\ Tetapkan

\Centerline{$H_n=\{x:x\in F,\,|f_{F,n+1}(x)-f_{Fn}(x)|\ge 2^{-n}\}$.}

\noindent Lalu $\mu H_n\le 2^{-n}$ untuk setiap $n$, jadi

\Centerline{$\mu(\bigcap_{n\in\Bbb N}\bigcup_{m\ge n}H_m)
\le\inf_{n\in\Bbb
N}\sum_{m=n}^{\infty}2^{-m}
=0$.}

\noindent Jika $x\in F\setminus\bigcap_{n\in\Bbb N}\bigcup_{m\ge n}H_m$, maka
terdapat $k$ sehingga $x\in F\setminus\bigcup_{m\ge k}H_m$, sehingga
$|f_{F,m+1}(x)-f_{Fm}(x)|\le 2^{-m}$ untuk setiap $m\ge k$, dan
$\sequencen{f_{Fn}(x)}$ adalah Cauchy, oleh karena itu konvergen.\ \Qed

Tetapkan $f_F(x)=\lim_{n\to\infty}f_{Fn}(x)$ untuk setiap $x\in F$ yang
batasnya ada di $\Bbb R$; dengan demikian $f_F$ terukur dan terdefinisi
hampir di mana-mana di $F$.

Jika $E$, $F$ adalah dua himpunan berukuran berhingga dan $E\subseteq F$, maka
$\bar\rho_E(u_{En},u_{Fn})\le 2\cdot 4^{-n}$ untuk setiap $n$.   \Prf\
$A_n(E)$ dan $A_n(F)$ keduanya milik $\Cal F$, sehingga keduanya beririsan;
ambil titik $w$ dalam irisan tersebut. Maka

$$\eqalign{\bar\rho_E(u_{En},u_{Fn})
&\le\bar\rho_E(u_{En},w)+\bar\rho_E(w,u_{Fn})\cr
&\le\bar\rho_E(u_{En},w)+\bar\rho_F(w,u_{Fn})
\le 4^{-n}+4^{-n}.  \text{
\Qed}\cr}$$



\noindent Akibatnya

\Centerline{$\mu\{x:x\in E,\,|f_{Fn}(x)-f_{En}(x)|\ge 2^{-n}\}
\le
2^n\bar\rho_E(u_{Fn},u_{En})\le 2^{-n+1}$}

\noindent untuk setiap $n$, dan $\lim_{n\to\infty}f_{Fn}-f_{En}=0$ hampir di
mana-mana pada $E$; sehingga $f_E=f_F$ hampir di mana-mana pada $E$.

Akibatnya, jika $E$ dan $F$ adalah dua himpunan berukuran berhingga,
$f_E=f_F$ hampir di mana-mana pada $E\cap F$, karena pada $E\cap F$ keduanya
sama hampir di mana-mana dengan $f_{E\cup F}$.

Karena $\mu$ dapat dilokalkan, terdapat $f\in\eusm L^0$ sehingga $f=f_E$
hampir di mana-mana pada $E$ untuk setiap himpunan terukur $E$ berukuran
berhingga (213N). Pertimbangkan $u=f^{\ssbullet}\in L^0$. Untuk setiap
himpunan $E$ berukuran berhingga,

\Centerline{$\bar\rho_E(u,u_{En})
=\int_E\min(1,|f(x)-f_{En}(x)|)dx
=\int_E\min(1,|f_E(x)-f_{En}(x)|)dx\to 0$}

\noindent ketika $n\to\infty$, menurut Teorema Konvergensi Terdominasi
Lebesgue. Selanjutnya,

$$\eqalign{\inf_{A\in\Cal F}\sup_{v\in A}\bar\rho_E(v,u)
&\le\inf_{n\in\Bbb
N}\sup_{v\in A_{En}}\bar\rho_E(v,u)\cr
&\le\inf_{n\in\Bbb N}\sup_{v\in A_{En}}
(\bar\rho_E(v,u_{En})+\bar\rho_E(u,u_{En}))\cr
&\le\inf_{n\in\Bbb
N}(4^{-n}+\bar\rho_E(u,u_{En}))
=0.\cr}$$

\noindent Karena $E$ sembarang, $\Cal F\to u$. Karena $\Cal F$ sembarang,
$L^0$ lengkap terhadap $\frak T$.

\medskip

\quad{\bf (ii)} Misalkan $L^0$ lengkap terhadap $\frak T$, dan
ambil sebarang keluarga himpunan $\Cal E$ dalam $\Sigma$. Misalkan
$\Cal E'$ adalah

\Centerline{$\{\bigcup\Cal E_0:\Cal E_0$ adalah subhimpunan hingga dari $\Cal
E\}$.}

\noindent Maka gabungan setiap dua anggota $\Cal E'$ termasuk dalam
$\Cal E'$. Misalkan $\Cal F$ adalah himpunan

\Centerline{$\{A:A\subseteq L^0$, $A\supseteq A_E$ untuk beberapa $E\in\Cal
E'\}$,}

\noindent di mana untuk $E\in\Cal E'$ saya tulis

\Centerline{$A_E=\{\chi F^{\ssbullet}:F\in\Cal E'$, $F\supseteq E\}$.}

\noindent Maka $\Cal F$ merupakan filter pada $L^0$, karena $A_E\cap
A_F=A_{E\cup
F}$ untuk semua $E$, $F\in\Cal E'$.

Sesungguhnya $\Cal F$ merupakan filter Cauchy. \Prf\ Misalkan $H$ suatu
himpunan berukuran berhingga dan $\epsilon>0$. Tetapkan $\gamma=\sup_{E\in\Cal
E'}\mu(H\cap E)$ dan ambil $E\in\Cal E'$ sehingga $\mu(H\cap E)\ge
\gamma-\epsilon$.   Pertimbangkan $A_E\in \Cal F$.   Jika $F$, $G\in\Cal E'$
dan $F\supseteq E$, $G\supseteq E$ maka

$$\eqalign{\bar\rho_H(\chi F^{\ssbullet},\chi G^{\ssbullet})
&=\mu(H\cap(F\symmdiff G))
=\mu(H\cap(F\cup G))-\mu(H\cap F\cap G)\cr
&\le\gamma-\mu(H\cap E)
\le\epsilon.\cr}$$

\noindent Jadi $\bar\rho_H(u,v)\le\epsilon$ untuk semua $u$, $v\in A_E$.
Karena $H$ dan $\epsilon$ sembarang, $\Cal F$ adalah Cauchy.\ \Qed

Karena $L^0$ lengkap terhadap $\frak T$, filter $\Cal F$ mempunyai limit
$w$. Nyatakan $w$ sebagai $h^{\ssbullet}$, dengan
$h:X\to\Bbb R$ terukur, dan pertimbangkan $G=\{x:h(x)>\bover12\}$.

\vthsp\Quer\ Jika $E\in \Cal E$ dan $\mu(E\setminus G)>0$, misalkan
$F\subseteq E\setminus G$ suatu himpunan berukuran positif dan berhingga.
Maka $\{u:\bar\rho_F(u,w)<\bover12\mu F\}$ termasuk dalam $\Cal F$,
sehingga beririsan dengan $A_E$; pilih $H\in\Cal E'$ sedemikian sehingga
$E\subseteq H$ dan
$\bar\rho_F(\chi H^{\ssbullet},w)<\bover12\mu F$. Maka

\Centerline{$\int_F\min(1,|1-h(x)|)
=\bar\rho_F(\chi
H^{\ssbullet},w)<\Bover12\mu F$;}

\noindent tetapi karena $F\cap G=\emptyset$, $1-h(x)\ge\bover12$ untuk setiap
$x\in F$, jadi ini tidak mungkin.\ \Bang

Jadi $E\setminus G$ terabaikan untuk setiap $E\in\Cal E$.

Selanjutnya, misalkan $H\in\Sigma$ dan $\mu(G\setminus H)>0$. Maka terdapat
$E\in\Cal E$ sehingga $\mu(E\setminus H)>0$.   \Prf\ Misalkan
$F\subseteq G\setminus H$ suatu himpunan berukuran positif dan berhingga.
Pilih $u\in A_{\emptyset}$ sedemikian sehingga
$\bar\rho_F(u,w)<\bover12\mu F$. Maka $u$ berbentuk
$\chi C^{\ssbullet}$ untuk beberapa $C\in\Cal E'$, dan

\Centerline{$\int_F\min(1,|\chi C(x) - h(x)|)<\Bover12\mu F$.}

\noindent Karena $h(x)\ge\bover12$ untuk setiap $x\in F$, $\mu(C\cap F)>0$.
Tetapi $C$ merupakan gabungan hingga dari anggota $\Cal E$, jadi terdapat
$E\in\Cal E$
sehingga $\mu(E\cap F)>0$, dan sekarang $\mu(E\setminus
H)>0$.\ \Qed

Karena $H$ sembarang, $G$ merupakan supremum esensial dari $\Cal E$ dalam
$\Sigma$. Karena $\Cal E$ sembarang, $(X,\Sigma,\mu)$ dapat
dilokalkan. }%akhir pembuktian 245E

\leader{245F}{Deskripsi alternatif topologi konvergensi dalam ukuran} Mari
kita kembali ke sebarang ruang ukur $(X,\Sigma,\mu)$.

\header{245Fa}{\bf (a)} Untuk setiap $F\in\Sigma$ berukuran berhingga dan
$\epsilon>0$, definisikan $\tau_{F\epsilon}:\eusm L^0\to\coint{0,\infty}$
melalui rumus

\Centerline{$\tau_{F\epsilon}(f)
=\mu^*\{x:x\in F\cap\dom
f,\,|f(x)|>\epsilon\}$}

\noindent untuk $f\in\eusm L^0$\cmmnt{, dengan $\mu^*$ menyatakan ukuran luar
yang ditentukan dari $\mu$ (132B)}. Jika $f$, $g\in\eusm L^0$
dan $f\eae g$, maka\cmmnt{

\Centerline{$\{x:x\in F\cap\dom f,\,|f(x)|>\epsilon\}
\symmdiff\{x:x\in
F\cap\dom g,\,|g(x)|>\epsilon\}$}

\noindent terabaikan, jadi} $\tau_{F\epsilon}(f)=\tau_{F\epsilon}(g)$;
oleh karena itu diperoleh fungsional dari $L^0$ ke $\coint{0,\infty}$
yang diberikan oleh

\Centerline{$\bar\tau_{F\epsilon}(u)=\tau_{F\epsilon}(f)$}

\noindent setiap kali $f\in\eusm L^0$ dan $u=f^{\ssbullet}\in L^0$.

\header{245Fb}{\bf (b)} \cmmnt{Fungsional $\tau_{F\epsilon}$ tidak subaditif
(kecuali dalam kasus-kasus sepele), sehingga tidak menentukan pseudometrik
pada $\eusm L^0$ maupun $L^0$. Namun kita dapat mengatakan bahwa, untuk
$f\in\eusm L^0$,

\Centerline{$\tau_F(f)\le\epsilon\min(1,\epsilon)
\Longrightarrow\tau_{F\epsilon}(f)\le\epsilon
\Longrightarrow\tau_F(f)\le\epsilon(1+\mu F)$.}

\noindent\prooflet{(Alasannya, jika $E\subseteq\dom f$ adalah
himpunan terukur yang koterabaikan dan $f\restr E$ terukur, maka

\Centerline{$\tau_F(f)=\int_{E\cap F}\min(|f(x)|,1)dx$,
\quad$\tau_{F\epsilon}(f)=\mu\{x:x\in E\cap F,\,|f(x)|>\epsilon\}$.)} }

Ini berarti bahwa himpunan $G\subseteq\eusm L^0$ terbuka untuk topologi
konvergensi dalam ukuran jika dan hanya jika, untuk setiap $f\in G$, terdapat
himpunan $F$ berukuran berhingga dan $\epsilon$, $\delta>0$ sedemikian rupa
sehingga

\Centerline{$\tau_{F\epsilon}(g-f)\le\delta\Longrightarrow g\in G$.}

\noindent Tentu saja $\tau_{F\delta}(f)\ge\tau_{F\epsilon}(f)$ setiap kali
$\delta\le\epsilon$, sehingga pernyataan itu dapat pula dirumuskan sebagai berikut: }%akhir komentar
$G\subseteq\eusm L^0$ terbuka untuk topologi konvergensi dalam ukuran jika
dan hanya jika, untuk setiap $f\in G$, terdapat himpunan $F$ berukuran berhingga dan
$\epsilon>0$ sedemikian rupa sehingga

\Centerline{$\tau_{F\epsilon}(g-f)\le\epsilon\Longrightarrow g\in G$.}

\noindent\cmmnt{Demikian pula, }$G\subseteq L^0$ terbuka untuk topologi
konvergensi dalam ukuran pada $L^0$ jika dan hanya jika, untuk setiap
$u\in G$, terdapat himpunan $F$ berukuran berhingga dan $\epsilon>0$ sedemikian rupa
sehingga

\Centerline{$\bar\tau_{F\epsilon}(v-u)\le\epsilon
\Longrightarrow v\in G$.}

\header{245Fc}{\bf (c)} Langsung diperoleh bahwa barisan $\sequencen{f_n}$
dalam $\eusm L^0=\eusm L^0(\mu)$ konvergen dalam ukuran ke $f\in\eusm L^0$
jika dan hanya jika

\Centerline{$\lim_{n\to\infty}
  \mu^*\{x:x\in F\cap\dom f\cap\dom
f_n,\,|f_n(x)-f(x)|>\epsilon\}
=0$}

\noindent setiap kali $F\in\Sigma$, $\mu F<\infty$ dan $\epsilon>0$. Demikian
pula, barisan $\sequencen{u_n}$ dalam $L^0$ konvergen dalam ukuran ke $u$ jika
dan hanya jika
$\lim_{n\to\infty}\bar\tau_{F\epsilon}(u-u_n)=0$ setiap kali $\mu
F<\infty$
dan $\epsilon>0$.

\header{245Fd}{\bf (d)} Secara khusus, jika $(X,\Sigma,\mu)$ berhingga total,
$\sequencen{f_n}\to f$ dalam $\eusm L^0$ jika dan hanya jika

\Centerline{$\lim_{n\to\infty}
  \mu^*\{x:x\in \dom f\cap\dom
f_n,\,|f(x)-f_n(x)|>\epsilon\}
=0$}

\noindent untuk setiap $\epsilon>0$, dan $\sequencen{u_n}\to u$ dalam $L^0$
jika dan hanya jika

\Centerline{$\lim_{n\to\infty}\bar\tau_{X\epsilon}(u-u_n)=0$}

\noindent untuk setiap $\epsilon>0$.

\leader{245G}{Penyematan $L^p$ dalam $L^0$: Proposisi} Misalkan
$(X,\Sigma,\mu)$ sebarang ruang ukur. Maka, untuk setiap $p\in[1,\infty]$,
penyematan $L^p=L^p(\mu)$ dalam $L^0=L^0(\mu)$ kontinu untuk topologi norma
$L^p$ dan topologi konvergensi dalam ukuran pada $L^0$.

\proof{ Misalkan $u$, $v\in L^p$ dan $\mu F<\infty$. Maka $(\chi
F)^{\ssbullet}$ merupakan anggota $L^q$, dengan $q=p/(p-1)$ (ambil $q=1$ jika
$p=\infty$, $q=\infty$ jika $p=1$ seperti biasa), dan

\Centerline{$\bar\rho_F(u,v)\le\int|u-v|\times(\chi F)^{\ssbullet}
\le\|u-v\|_p\|\chi F^{\ssbullet}\|_q$}

\noindent (244Eb). Menurut 2A3H, ini cukup untuk memastikan bahwa penyematan
$L^p\embedsinto L^0$ kontinu. }% akhir bukti 245G

\leader{245H}{}\cmmnt{ Kasus $L^1$ begitu penting sehingga saya membahasnya
lebih terperinci.

\medskip

\noindent}{\bf Proposisi} Misalkan $(X,\Sigma,\mu)$ adalah ruang ukur.

(a)(i) Jika $f\in\eusm L^1=\eusm L^1(\mu)$ dan $\epsilon>0$, terdapat
$\delta>0$ dan himpunan $F\in\Sigma$ berukuran berhingga sehingga
$\int|f-g|\le\epsilon$ setiap kali $g\in\eusm L^1$, $\int|g|\le\int|f|+\delta$
dan $\rho_F(f,g)\le\delta$.

\quad(ii) Untuk setiap barisan $\sequencen{f_n}$ dalam $\eusm L^1$ dan $f\in\eusm
L^1$, $\lim_{n\to\infty}\int|f-f_n|=0$ jika dan hanya jika $\sequencen{f_n}\to f$ dalam ukuran
dan $\limsup_{n\to\infty}\int|f_n|\le\int|f|$.

(b)(i) Jika $u\in L^1=L^1(\mu)$ dan $\epsilon>0$, terdapat $\delta>0$ dan
himpunan $F\in\Sigma$ berukuran berhingga sehingga $\|u-v\|_1\le\epsilon$
setiap kali $v\in L^1$, $\|v\|_1\le\|u\|_1+\delta$ dan
$\bar\rho_F(u,v)\le\delta$.

\quad(ii) Untuk setiap barisan $\sequencen{u_n}$ dalam $L^1$ dan setiap $u\in
L^1$, $\sequencen{u_n}\to u$ untuk $\|\,\|_1$ jika dan hanya jika $\sequencen{u_n}\to u$ dalam
ukuran dan $\limsup_{n\to\infty}\|u_n\|_1\le\|u\|_1$.

\proof{{\bf (a)(i)} Menurut 225A, terdapat himpunan $F$ berukuran berhingga
dan $\eta>0$ sehingga $\int_E|f|\le\bover15\epsilon$ setiap kali $\mu(E\cap
F)\le\eta$. Ambil $\delta>0$ sehingga $\delta(\epsilon+5\mu
F)\le\epsilon\eta$ dan $\delta\le\bover15\epsilon$. Lalu jika
$\int|g|\le\int|f|+\delta$ dan $\rho_F(f,g)\le\delta$, misalkan
$G\subseteq\dom f\cap\dom g$ adalah himpunan terukur yang koterabaikan
sehingga $f\restr G$ dan $g\restr G$ keduanya terukur. Tetapkan

\Centerline{$E=\{x:x\in F\cap G,\,
|f(x)-g(x)|\ge\Bover{\epsilon}{\epsilon+5\mu F}\}$;}

\noindent lalu

\Centerline{$\delta\ge\rho_F(f,g)
\ge\Bover{\epsilon}{\epsilon+5\mu F}\mu E$,}

\noindent jadi $\mu E\le\eta$.   Tetapkan $H=F\setminus E$, sehingga
$\mu(F\setminus H)\le\eta$ dan $\int_{X\setminus
H}|f|\le\bover15{\epsilon}$.
Di sisi lain, untuk hampir setiap $x\in
H$,
$|f(x)-g(x)|\le\bover{\epsilon}{\epsilon+5\mu F}$, jadi
$\int_H|f-g|\le\bover15\epsilon$ dan

\Centerline{$\int_{H}|g|\ge\int_{H}|f|-\Bover15\epsilon
\ge\int|f|-\int_{X\setminus H}|f|-\Bover15\epsilon
\ge\int|f|-\Bover25\epsilon$.}

\noindent Karena $\int|g|\le\int|f|+\bover15\epsilon$, $\int_{X\setminus
H}|g|\le\bover35{\epsilon}$. Dengan demikian,

\Centerline{$\int|g-f|
\le\int_{X\setminus H}|g|+\int_{X\setminus
H}|f|+\int_H|g-f|
\le\Bover35\epsilon+\Bover15\epsilon+\Bover15\epsilon
=\epsilon$,}

\noindent sesuai kebutuhan.

\medskip

\quad{\bf (ii)} Jika $\lim_{n\to\infty}\int|f-f_n|=0$, yakni
$\lim_{n\to\infty}f_n^{\ssbullet}=f^{\ssbullet}$ di $L^1$, maka menurut
245G diperoleh $\sequencen{f_n^{\ssbullet}}\to f^{\ssbullet}$ dalam
$L^0$, yaitu, $\sequencen{f_n}\to f$ untuk topologi konvergensi dalam ukuran.
Selain itu, tentu saja, $\lim_{n\to\infty}\int|f_n|=\int|f|$.

Sebaliknya, jika $\limsup_{n\to\infty}\int|f_n|\le\int|f|$ dan
$\sequencen{f_n}\to f$ untuk topologi konvergensi dalam ukuran, maka untuk setiap
$\delta>0$ dan $\mu F<\infty$ terdapat $m\in\Bbb N$ sedemikian sehingga
$\int|f_n|\le\int|f|+\delta$, $\rho_F(f,f_n)\le\delta$ untuk setiap $n\ge m$;
jadi (i) menunjukkan bahwa $\lim_{n\to\infty}\int|f_n-f|=0$.

\medskip

{\bf (b)} Bagian ini langsung mengikuti jika kita menyatakan $u$ sebagai
$f^{\ssbullet}$, $v$ sebagai $g^{\ssbullet}$ dan $u_n$ sebagai
$f_n^{\ssbullet}$. }%akhir dari bukti 245H

\cmmnt{ \vleader{48pt}{245I}{Catatan (a)} Fenomena ini begitu penting sehingga
layak ditinjau melalui beberapa contoh dasar.

(i) Jika $\mu$ adalah ukuran hitung pada $\Bbb N$, dan kita menetapkan
$f_n(n)=1$, $f_n(i)=0$ untuk $i\ne n$, maka $\sequencen{f_n}\to 0$ dalam
ukuran, sementara $\int|f_n|=1$ untuk setiap $n$.

(ii) Jika $\mu$ adalah ukuran Lebesgue pada $[0,1]$, dan kita menetapkan
$f_n(x)=2^n$ untuk $0<x\le 2^{-n}$, $0$ untuk $x$ lainnya, maka sekali lagi
$\sequencen{f_n}\to 0$ dalam ukuran, sementara $\int|f_n|=1$ untuk setiap $n$.

(iii) Dalam 245Cc kita mempunyai barisan lain $\sequencen{f_n}$ yang konvergen
ke $0$ dalam ukuran, sementara $\int|f_n|=1$ untuk setiap $n$. Dalam semua
kasus ini (seperti yang disyaratkan oleh Lemma Fatou, setidaknya pada (i) dan
(ii)) kita mempunyai $\int|f|\le\liminf_{n\to\infty}\int|f_n|$. (Proposisi
berikut menunjukkan bahwa ini berlaku untuk setiap barisan yang konvergen dalam
ukuran.)

Kesamaan ketiga contoh ini ialah bahwa $f_n$ seolah-olah lolos menuju tak
hingga, baik secara lateral (dalam (i)) maupun secara
vertikal (dalam (iii)) atau keduanya (dalam (ii)). Perlu dicatat bahwa dalam
ketiga contoh tersebut kita dapat menetapkan $f'_n=2^nf_n$ untuk memperoleh
sebuah barisan yang masih konvergen ke $0$ dalam ukuran, tetapi dengan
$\lim_{n\to\infty}\int|f'_n|=\infty$.

\spheader 245Ib Dalam 245H, saya telah menggunakan rumusan eksplisit
`$\lim_{n\to\infty}\int|f_n-f|=0$' (untuk barisan fungsi), `$\sequencen{u_n}\to
u$ untuk $\|\,\|_1$' (untuk barisan dalam $L^1$). Hal ini sering dirumuskan
dengan menyebut $\sequencen{f_n}$ dan $\sequencen{u_n}$ masing-masing {\bf
konvergen dalam rata-rata} ke $f$ dan $u$. }%akhir komentar

\leader{245J}{}\cmmnt{ Untuk ruang semihingga kita mempunyai hubungan lebih
lanjut.

\medskip

\noindent}{\bf Proposisi} Misalkan $(X,\Sigma,\mu)$ adalah ruang ukur semihingga.
Tulis $\eusm L^0=\eusm L^0(\mu)$, dan seterusnya.

(a)(i) Untuk setiap $a\ge 0$, himpunan $\{f:f\in\eusm L^1,\,\int|f|\le a\}$
tertutup di $\eusm L^0$ untuk topologi konvergensi dalam ukuran.

\quad(ii) Jika $\sequencen{f_n}$ suatu barisan dalam $\eusm L^1$ yang
konvergen dalam ukuran ke $f\in\eusm L^0$, dan
$\liminf_{n\to\infty}\int|f_n|<\infty$, maka $f$ terintegralkan dan
$\int|f|\le\liminf_{n\to\infty}\int|f_n|$.

(b)(i) Untuk setiap $a\ge 0$, himpunan $\{u:u\in L^1,\,\|u\|_1\le a\}$
tertutup dalam $L^0$ untuk topologi konvergensi dalam ukuran.

\quad(ii) Jika $\sequencen{u_n}$ suatu barisan dalam $L^1$ yang konvergen
dalam ukuran ke $u\in L^0$, dan $\liminf_{n\to\infty}\|u_n\|_1<\infty$, maka
$u\in L^1$ dan $\|u\|_1\le\liminf_{n\to\infty}\|u_n\|_1$.

\proof{{\bf (a)(i)} Tetapkan $A=\{f:f\in\eusm L^1,\,\int|f|\le a\}$, dan
ambil $g$ sebarang anggota penutupan $A$ dalam $\eusm L^0$. Misalkan $h$
fungsi sederhana sedemikian sehingga $0\le h\leae|g|$, dan
$\epsilon>0$.   Jika $h=0$ maka tentu saja $\int h\le a$.   Jika tidak, dengan
menetapkan $F=\{x:h(x)>0\}$ dan $M=\sup_{x\in X}h(x)$, terdapat $f\in A$ sedemikian
sehingga $\mu^*\{x:x\in F\cap\dom f\cap\dom g,\,
|f(x)-g(x)|\ge\epsilon\}\le\epsilon$ (245F); pilih himpunan terukur
$E\supseteq\{x:x\in F\cap\dom f\cap\dom g,\,|f(x)-g(x)|\ge\epsilon\}$
yang berukuran paling besar $\epsilon$. Maka $h\leae|f|+\epsilon\chi F+M\chi
E$, sehingga $\int h\le a+\epsilon(M+\mu F)$. Karena $\epsilon$ sembarang,
$\int h\le a$. Karena $\mu$ semihingga, hal ini
cukup untuk memastikan bahwa $g$ terintegralkan
dan bahwa $\int|g|\le a$ (213B), yakni $g\in A$. Karena $g$ dipilih
sembarang, $A$ tertutup.

\medskip

\quad{\bf (ii)} Jika $\sequencen{f_n}$ konvergen dalam ukuran ke $f$,
dan $\liminf_{n\to\infty}\int|f_n|=a$, maka untuk setiap $\epsilon>0$ terdapat
subbarisan $\sequence{k}{f_{n(k)}}$ sedemikian rupa sehingga $\int|f_{n(k)}|\le
a+\epsilon$ untuk setiap $k$; karena $\sequence{k}{f_{n(k)}}$ masih konvergen
dalam ukuran ke $f$, $\int|f|\le a+\epsilon$. Karena $\epsilon$ sembarang,
$\int|f|\le a$.

\medskip

{\bf (b)} Seperti pada 245H, ini hanyalah terjemahan dari bagian (a) ke dalam
bahasa $L^1$ dan $L^0$. }%akhir dari bukti 245J


\leader{245K}{}\cmmnt{Untuk ruang ukur $\sigma$-hingga, topologi konvergensi
dalam ukuran pada $L^0$ dapat dimetriskan, sehingga dapat dijelaskan secara
efektif melalui barisan konvergen; karena itu, dalam kasus ini penting
untuk mempunyai karakterisasi yang tepat bagi konvergensi barisan
dalam ukuran.

\medskip

\noindent}{\bf Proposisi} Misalkan $(X,\Sigma,\mu)$ ruang ukur
$\sigma$-hingga. Maka

(a) suatu barisan $\sequencen{f_n}$ dalam $\eusm L^0$ konvergen dalam ukuran
ke $f\in\eusm L^0$ jika dan hanya jika setiap subbarisan dari $\sequencen{f_n}$
mempunyai subsubbarisan yang konvergen ke $f$ hampir di mana-mana;

(b) suatu barisan $\sequencen{u_n}$ dalam $L^0$ konvergen dalam ukuran ke
$u\in
L^0$ jika dan hanya jika setiap subbarisan dari $\sequencen{u_n}$
mempunyai subsubbarisan yang ordo*-konvergen ke $u$.

\proof{{\bf (a)(i)} Andaikan $\sequencen{f_n}\to f$, yaitu
$\lim_{n\to\infty}\int|f-f_n|\wedge \chi F=0$ untuk setiap himpunan $F$ dengan
ukuran berhingga. Misalkan $\sequencen{E_k}$ barisan tak menurun dari himpunan
berukuran berhingga yang menutupi $X$, dan misalkan
$\sequence{k}{n(k)}$ barisan naik tegas dalam $\Bbb N$
sedemikian rupa sehingga $\int|f-f_{n(k)}|\wedge\chi E_k\le 4^{-k}$ untuk
setiap $k\in\Bbb N$. Maka $\sum_{k=0}^{\infty}|f-f_{n(k)}|\wedge\chi E_k$
adalah hingga hampir di mana-mana (242E); tetapi
$\lim_{k\to\infty}f_{n(k)}(x)=f(x)$ setiap kali
$\sum_{k=0}^{\infty}\min(1,|f(x)-f_{n(k)}(x)|)<\infty$, jadi
$\sequence{k}{f_{n(k)}}\to f$ hampir di mana-mana.

\medskip

\quad{\bf (ii)} Hal yang sama berlaku bagi setiap subbarisan dari
$\sequencen{f_n}$, sehingga setiap subbarisan dari $\sequencen{f_n}$ mempunyai
subsubbarisan yang konvergen ke $f$ hampir di mana-mana.

\medskip

\quad{\bf (iii)} Misalkan sekarang bahwa $\sequencen{f_n}$ tidak konvergen ke
$f$. Maka ada himpunan terbuka $U$ yang memuat $f$ sehingga
$\{n:f_n\notin U\}$ tak hingga, yaitu, $\sequencen{f_n}$ mempunyai subbarisan
$\sequencen{f'_n}$ dengan $f'_n\notin U$ untuk setiap $n$. Namun
tidak ada subsubbarisan dari $\sequencen{f'_n}$ yang konvergen ke $f$ dalam
ukuran, sehingga tidak ada subsubbarisan seperti itu yang dapat konvergen
hampir di mana-mana, menurut 245Ca.

\medskip

{\bf (b)} Ini langsung mengikuti (a) jika $u$ dinyatakan sebagai
$f^{\ssbullet}$ dan $u_n$ sebagai $f_n^{\ssbullet}$. }%akhir pembuktian 245K


\leader{245L}{Akibat} Misalkan $(X,\Sigma,\mu)$ ruang ukur $\sigma$-hingga.

(a) Suatu subhimpunan $A$ dari $\eusm L^0=\eusm L^0(\mu)$ tertutup untuk
topologi konvergensi dalam ukuran jika dan hanya jika $f\in A$ setiap kali
$f\in\eusm L^0$ dan terdapat barisan $\sequencen{f_n}$ dalam $A$ sedemikian
sehingga $f\eae\lim_{n\to\infty}f_n$.

(b) Suatu subhimpunan $A$ dari $L^0=L^0(\mu)$ tertutup untuk topologi
konvergensi dalam ukuran jika dan hanya jika $u\in A$ setiap kali $u\in L^0$
dan terdapat suatu barisan $\sequencen{u_n}$ dalam $A$ yang ordo*-konvergen
ke $u$.

\proof{{\bf (a)(i)} Jika $A$ tertutup untuk topologi konvergensi dalam ukuran,
dan $\sequencen{f_n}$ merupakan barisan dalam $A$ yang konvergen ke $f$
hampir di mana-mana, maka $\sequencen{f_n}$ konvergen ke $f$ dalam ukuran,
jadi tentu saja $f\in A$ (karena jika tidak, maka hampir semua
$f_n$ harus termasuk dalam himpunan terbuka $\eusm
L^0\setminus A$).

\medskip

\quad{\bf (ii)} Jika $A$ tidak tertutup, ada $f\in\overline{A}\setminus A$.
Topologi dapat ditentukan oleh metrik $\rho$ (245Eb), dan kita dapat
memilih barisan $\sequencen{f_n}$ dalam $A$ sedemikian sehingga
$\rho(f_n,f)\le 2^{-n}$ untuk setiap $n$, sehingga $\sequencen{f_n}\to f$
dalam ukuran. Menurut 245K, $\sequencen{f_n}$ mempunyai subbarisan
$\sequencen{f'_n}$ yang konvergen hampir di mana-mana ke $f$, dan ini membuktikan bahwa $A$
gagal memenuhi kondisi.

\medskip

{\bf (b)} Ini langsung mengikuti fakta bahwa $A\subseteq L^0$ tertutup jika dan hanya
jika $\{f:f^{\ssbullet}\in A\}$ tertutup di $\eusm L^0$. }%akhir pembuktian 245L

\leader{245M}{$L^0$ kompleks} Dalam 241J saya secara singkat
membahas adaptasi yang diperlukan untuk membangun ruang linier kompleks
$L^0_{\Bbb C}$. Rumus-rumus dari 245A dapat digunakan tanpa perubahan untuk
mendefinisikan topologi konvergensi dalam ukuran pada $\eusm L^0_{\Bbb C}$ dan
$L^0_{\Bbb C}$. Tampaknya setiap pernyataan dalam 245B-245L tetap berlaku jika kita
mengganti setiap $L^0$ atau $\eusm L^0$ dengan $L^0_{\Bbb C}$ atau $\eusm
L^0_{\Bbb C}$. \cmmnt{Sebagai alternatif, untuk menghubungkan bentuk `riil'
dan `kompleks' dari 245E, misalnya, perhatikan bahwa

$$\eqalign{\max(\rho_F(\Real(u),\Real(v)),\rho_F(\Imag(u),\Imag(v)))
&\le\rho_F(u,v)\cr
&\le\rho_F(\Real(u),\Real(v))+\rho_F(\Imag(u),\Imag(v))\cr}$$

\noindent untuk semua $u$, $v\in L^0$ dan semua himpunan $F$ berukuran
berhingga, $L^0_{\Bbb C}$ dapat diidentifikasi, sebagai ruang seragam, dengan
$L^0\times
L^0$, sehingga bersifat Hausdorff, dapat dimetriskan, atau lengkap jika
dan hanya jika $L^0$ demikian.}

\exercises{ \leader{245X}{Latihan dasar $\pmb{>}$(a)}
%\spheader 245Xa
Misalkan $X$ sebarang himpunan, dan $\mu$ ukuran hitung pada $X$.
Tunjukkan bahwa topologi konvergensi dalam ukuran pada $\eusm L^0(\mu)=\Bbb
R^X$ hanyalah topologi produk pada $\Bbb R^X$ jika dianggap sebagai produk
dari salinan $\Bbb R$.
%245A

\sqheader 245Xb Misalkan $(X,\Sigma,\mu)$ adalah ruang ukur sembarang, dan
$(X,\hat\Sigma,\hat\mu)$ adalah kelengkapannya. Tunjukkan bahwa topologi
konvergensi dalam ukuran pada $\eusm L^0(\mu)=\eusm L^0(\hat\mu)$ (241Xb),
yang sesuai dengan keluarga $\{\rho_F:F\in\Sigma,\,\mu F<\infty\}$,
$\{\rho_F:F\in\hat\Sigma,\,\hat\mu F<\infty\}$ adalah sama.
%245A

\sqheader 245Xc Misalkan $(X,\Sigma,\mu)$ adalah sebarang ruang ukur; tetapkan
$L^0=L^0(\mu)$. Misalkan $u$, $u_n\in L^0$ untuk $n\in\Bbb N$. Tunjukkan bahwa
pernyataan-pernyataan berikut ekuivalen:

\quad(i) $\sequencen{u_n}$ ordo*-konvergen ke $u$ dalam pengertian 245C;

\quad (ii) terdapat fungsi terukur $f$, $f_n:X\to\Bbb R$ sehingga
$f^{\ssbullet}=u$, $f_n^{\ssbullet}=u_n$ untuk setiap $n\in\Bbb N$, dan
$f(x)=\lim_{n\to\infty}f_n(x)$ untuk setiap $x\in X$;

\quad (iii) $u=\inf_{n\in\Bbb N}\sup_{m\ge n}u_m=\sup_{n\in\Bbb N}\inf_{m\ge
n}u_m$, infima dan suprema diambil dalam $L^0$;

\quad (iv) $\inf_{n\in\Bbb N}\sup_{m\ge n}|u-u_m|=0$ dalam $L^0$;

\quad (v) terdapat barisan tak meningkat $\sequencen{v_n}$ dalam $L^0$
sehingga $\inf_{n\in\Bbb N}v_n=0$ dalam $L^0$ dan $u-v_n\le u_n\le u+v_n$
untuk setiap $n\in\Bbb N$;

\quad (vi) terdapat barisan $\sequencen{v_n}$, $\sequencen{w_n}$ dalam $L^0$
sehingga $\sequencen{v_n}$ tak menurun, $\sequencen{w_n}$ tak meningkat,
$\sup_{n\in\Bbb N}v_n=u=\inf_{n\in\Bbb N}w_n$ dan $v_n\le u_n\le w_n$ untuk
setiap $n\in\Bbb N$.
%245C

\spheader 245Xd Misalkan $(X,\Sigma,\mu)$ ruang ukur semihingga.
Tunjukkan bahwa barisan $\sequencen{u_n}$ dalam $L^0=L^0(\mu)$
ordo*-konvergen ke $u\in L^0$ jika dan hanya jika $\{|u_n|:n\in\Bbb
N\}$ dibatasi di atas dalam $L^0$ dan $\sequencen{\sup_{m\ge n}|u_m-u|}\to 0$
untuk topologi konvergensi dalam ukuran.
%245C

\spheader 245Xe Tuliskan bukti bahwa $L^0(\mu)$ lengkap (sebagai ruang linier
topologis), yang disesuaikan dengan kasus-kasus khusus (i) $\mu X=1$ (ii)
$\mu$ $\sigma$-hingga, dengan memanfaatkan setiap penyederhanaan yang dapat
Anda temukan.
%245E

\spheader 245Xf Misalkan $(X,\Sigma,\mu)$ adalah ruang ukur dan $r\ge 1$;
misalkan $h:\Bbb R^r\to\Bbb R$ fungsi kontinu. (i) Andaikan untuk
$1\le k\le r$ kita diberikan barisan $\sequencen{f_{kn}}$ dalam $\eusm
L^0={\eusm L}^0(\mu)$ yang konvergen dalam ukuran ke $f_k\in{\eusm L}^0$.
Tunjukkan bahwa $\sequencen{h(f_{1n},\ldots,f_{rn})}$ konvergen dalam ukuran
ke $h(f_1,\ldots,f_r)$. (ii) Secara umum, tunjukkan bahwa
$(f_1,\ldots,f_r)\mapsto h(f_1,\ldots,f_r):(\eusm L^0)^r\to\eusm L^0$ adalah
kontinu untuk topologi konvergensi dalam ukuran.   (iii) Tunjukkan bahwa
fungsi yang bersesuaian $\bar h:(L^0)^r\to L^0$ (241Xh) adalah kontinu untuk
topologi konvergensi dalam ukuran.
%245F

\spheader 245Xg Misalkan $(X,\Sigma,\mu)$ adalah ruang ukur dan $u\in
L^1(\mu)$. Tunjukkan bahwa $v\mapsto\int u\times v:L^{\infty}\to\Bbb R$ adalah
kontinu untuk topologi konvergensi dalam ukuran pada bola satuan
$L^{\infty}$, tetapi pada umumnya tidak kontinu pada seluruh $L^{\infty}$.
%245G

\spheader 245Xh Misalkan $(X,\Sigma,\mu)$ merupakan ruang ukur dan $v$ adalah
anggota nonnegatif dari $L^1=L^1(\mu)$. Tunjukkan bahwa pada himpunan
$A=\{u:u\in L^1,\,|u|\le v\}$, topologi subruang (2A3C) yang diturunkan oleh
topologi norma pada $L^1$ dan topologi konvergensi dalam ukuran adalah sama.
\Hint{Jika diberikan $\epsilon>0$, ambil $F\in\Sigma$ berukuran berhingga dan
$M\ge 0$ sehingga $\int(|v|-M\chi F^{\ssbullet})^+\le\epsilon$. Tunjukkan
bahwa $\|u-u'\|_1\le\epsilon+M\bar\rho_F(u,u')$ untuk setiap $u$, $u'\in A$. }
%245H

\spheader 245Xi Misalkan $(X,\Sigma,\mu)$ adalah ruang ukur dan $\Cal F$
adalah filter pada $L^1=L^1(\mu)$ yang konvergen, untuk topologi konvergensi
dalam ukuran, ke $u\in L^1$. Tunjukkan bahwa $\Cal F\to u$ untuk topologi norma
$L^1$ jika dan hanya jika $\inf_{A\in\Cal F}\sup_{v\in
A}\|v\|_1\le\|u\|_1$.
%245H

\spheader 245Xj Misalkan $(X,\Sigma,\mu)$ adalah ruang ukur dan
$p\in\coint{1,\infty}$. Misalkan $\sequencen{u_n}$ barisan
dalam $L^p(\mu)$ yang konvergen untuk $\|\,\|_p$ ke $u\in L^p(\mu)$. Tunjukkan
bahwa $\sequencen{|u_n|^p}\to|u|^p$ untuk $\|\,\|_1$. \Hint{245G, 245Dd,
245H.}
%245H added 2002

\sqheader 245Xk Misalkan $(X,\Sigma,\mu)$ ruang ukur semihingga dan
$p\in[1,\infty]$, $a\ge 0$. Tunjukkan bahwa $\{u:u\in L^p(\mu),\,\|u\|_p\le
a\}$ tertutup di $L^0(\mu)$ untuk topologi konvergensi dalam ukuran.
%245J

\spheader 245Xl Misalkan $(X,\Sigma,\mu)$ adalah ruang ukur, dan
$\sequencen{u_n}$ barisan dalam $L^p=L^p(\mu)$, di mana $1\le
p<\infty$. Misalkan $u\in L^p$. Tunjukkan bahwa hal-hal berikut ini adalah
ekuivalen: (i) $u=\lim_{n\to\infty}u_n$ untuk topologi norma $L^p$ (ii)
$\sequencen{u_n}\to u$ untuk topologi konvergensi dalam ukuran dan
$\lim_{n\to\infty}\|u_n\|_p=\|u\|_p$ (iii) $\sequencen{u_n}\to u$ untuk
topologi konvergensi dalam ukuran dan
$\limsup_{n\to\infty}\|u_n\|_p\le\|u\|_p$.
%245Xk, 245J

\spheader 245Xm Misalkan $X$ suatu himpunan dan $\mu$, $\nu$ dua ukuran pada
$X$ dengan himpunan-himpunan terukur dan himpunan-himpunan terabaikan yang
sama. (i) Tunjukkan bahwa $\eusm L^0(\mu)=\eusm L^0(\nu)$ dan
$L^0(\mu)=L^0(\nu)$. (ii) Tunjukkan bahwa jika $\mu$ dan $\nu$ sama-sama
semihingga, maka keduanya menentukan topologi yang sama bagi
konvergensi dalam ukuran pada $\eusm L^0$ dan $L^0$. \Hint{Gunakan 215A
untuk menunjukkan bahwa jika $\mu E<\infty$ maka $\mu E=\sup\{\mu F:F\subseteq
E,\,\nu F<\infty\}$. }
%245L

\leader{245Y}{Latihan lebih lanjut (a)}
%\spheader 245Ya
Misalkan $(X,\Sigma,\mu)$ adalah ruang ukur dan lengkapi $\Sigma$ dengan topologi
yang dijelaskan dalam 232Ya. Tunjukkan bahwa $\chi:\Sigma\to\eusm L^0(\mu)$ adalah
homeomorfisme antara $\Sigma$ dan citranya $\chi[\Sigma]$ di $\eusm L^0$, jika
$\eusm L^0$ diberi topologi konvergensi dalam ukuran dan $\chi[\Sigma]$
diberi topologi subruang.
%245A

\spheader 245Yb Misalkan $(X,\Sigma,\mu)$ ruang ukur dan $Y$ sembarang
subhimpunan $X$; misalkan $\mu_Y$ ukuran subruang pada $Y$. Misalkan
$T:L^0(\mu)\to L^0(\mu_Y)$ peta kanonik yang
didefinisikan dengan menetapkan $T(f^{\ssbullet})=(f\restrp Y)^{\ssbullet}$
untuk setiap $f\in\eusm L^0(\mu)$ (241Yg). Tunjukkan bahwa $T$ bersifat
kontinu untuk topologi konvergensi dalam ukuran pada $L^0(\mu)$ dan
$L^0(\mu_Y)$.
%245A

\spheader 245Yc Misalkan $(X,\Sigma,\mu)$ adalah ruang ukur, dan $\tilde\mu$
adalah versi c.l.d.\  dari $\mu$. Tunjukkan bahwa peta $T:L^0(\mu)\to
L^0(\tilde\mu)$ yang diinduksi oleh inklusi $\eusm L^0(\mu)\subseteq\eusm
L^0(\tilde\mu)$ (241Yf) bersifat kontinu untuk topologi konvergensi dalam ukuran.
%245A

\spheader 245Yd Misalkan $(X,\Sigma,\mu)$ ruang ukur, dan berikan
$L^0=L^0(\mu)$ topologi konvergensi dalam ukuran.  Misalkan $A\subseteq L^0$
adalah himpunan tak kosong yang terarah ke bawah, dan andaikan bahwa
$\inf A=0$ dalam $L^0$. (i) Misalkan $F\in\Sigma$ himpunan berukuran
berhingga, dan definisikan $\bar\tau_F$ seperti dalam 245A; tunjukkan bahwa
$\inf_{u\in A}\bar\tau_F(u)=0$.  \Hint{Tetapkan $\gamma=\inf_{u\in
A}\bar\tau_F(u)$; temukan barisan menurun $\sequencen{u_n}$ dalam $A$
sehingga $\lim_{n\to\infty}\bar\tau_F(u_n)=\gamma$; tetapkan $v=(\chi
F)^{\ssbullet}\wedge\inf_{n\in\Bbb N}u_n$ dan tunjukkan bahwa $u\wedge v=v$
untuk setiap $u\in A$, sehingga $v=0$.}  (ii) Tunjukkan bahwa jika $U$ adalah
himpunan terbuka yang memuat $0$, terdapat $u\in A$ sehingga $v\in U$ setiap
kali $0\le v\le u$.
%245A

\spheader 245Ye Misalkan $(X,\Sigma,\mu)$ adalah ruang ukur. (i) Tunjukkan
bahwa untuk $u\in L^0=L^0(\mu)$ kita dapat mendefinisikan $\psi_a(u)$, untuk
$a\ge 0$, dengan menetapkan $\psi_a(u)=\mu\{x:|f(x)|>a\}$ setiap kali
$f:X\to\Bbb R$ adalah fungsi terukur dan $f^{\ssbullet}=u$. (ii) Definisikan
$\rho:L^0\times
L^0\to[0,1]$ dengan menetapkan
$\rho(u,v)=\min(\{1\}\cup\{a:a\ge 0,\,\psi_a(u-v)\le a\})$. Tunjukkan bahwa
$\rho$ adalah metrik pada $L^0$, bahwa $L^0$ lengkap terhadap $\rho$, dan
bahwa $+$, $-$, $\wedge$, $\vee:L^0\times L^0\to L^0$ kontinu terhadap $\rho$.
(iii) Tunjukkan bahwa $c\mapsto cu:\Bbb R\to L^0$ kontinu untuk setiap $u\in
L^0$ jika dan hanya jika $(X,\Sigma,\mu)$ berhingga total, dan bahwa dalam
kasus ini $\rho$ mendefinisikan topologi konvergensi dalam ukuran pada
$L^0$.
%245D

\spheader 245Yf Misalkan $(X,\Sigma,\mu)$ ruang ukur yang dapat dilokalkan dan
$A\subseteq L^0=L^0(\mu)$ himpunan tak kosong, terarah ke atas, dan terbatas
dalam arti ruang linier topologis (yakni,
sedemikian sehingga untuk setiap lingkungan $U$ dari $0$ di $L^0$ terdapat
$k\in\Bbb N$ sedemikian sehingga $A\subseteq kU$). Tunjukkan bahwa $A$
terbatas di atas di $L^0$, dan bahwa supremumnya termasuk dalam penutupannya.
%245E

\spheader 245Yg Misalkan $(X,\Sigma,\mu)$ adalah ruang ukur,
$p\in\coint{1,\infty}$ dan $v$ anggota nonnegatif dari $L^p=L^p(\mu)$.
Tunjukkan bahwa pada himpunan $A=\{u:u\in L^p,\,|u|\le v\}$ topologi subruang
yang diinduksi oleh topologi norma $L^p$ dan topologi konvergensi dalam
ukuran adalah sama.
%245F

\spheader 245Yh Misalkan $S$ adalah himpunan semua barisan $s:\Bbb
N\to\Bbb N$ sedemikian sehingga $\lim_{n\to\infty}s(n)=\infty$. Untuk setiap
$s\in S$, misalkan $(X_s,\Sigma_s,\mu_s)$ adalah $[0,1]$ dengan ukuran
Lebesgue, dan misalkan $(X,\Sigma,\mu)$ jumlah langsung dari
$\family{s}{S}{(X_s,\Sigma_s,\mu_s)}$ (214L). Untuk $s\in S$, $t\in[0,1]$,
$n\in\Bbb N$ tetapkan $h_n(s,t)=f_{s(n)}(t)$, di mana $\sequencen{f_n}$ adalah
barisan dari 245Cc. Tunjukkan bahwa $\sequencen{h_n}\to 0$ untuk topologi
konvergensi dalam ukuran pada $\eusm L^0(\mu)$, tetapi bahwa $\sequencen{h_n}$
tidak memiliki subbarisan yang konvergen ke $0$ hampir di mana-mana.
%245K

\spheader 245Yi Misalkan $X$ suatu himpunan. Andaikan diberikan suatu
relasi $\rightharpoonup$ antara barisan dalam $X$ dan anggota
$X$ sedemikian sehingga ($\alpha$) jika $x_n=x$ untuk setiap $n$ maka
$\sequencen{x_n}\rightharpoonup
x$ ($\beta$) $\sequencen{x'_n}\rightharpoonup
x$ setiap kali $\sequencen{x_n}\rightharpoonup x$ dan $\sequencen{x'_n}$ adalah
subbarisan dari $\sequencen{x_n}$. Tunjukkan bahwa terdapat topologi
$\frak T$ pada $X$ yang ditentukan sebagai berikut: suatu himpunan
bagian $G$ dari $X$ termasuk dalam $\frak T$ jika dan hanya jika, setiap kali
$\sequencen{x_n}$ barisan dalam $X$ dan $\sequencen{x_n}\rightharpoonup
x\in G$, setidaknya satu $x_n$ termasuk dalam $G$. Tunjukkan bahwa suatu barisan
$\sequencen{x_n}$ dalam $X$ $\frak T$-konvergen ke $x$ jika dan hanya jika
setiap subbarisan dari $\sequencen{x_n}$ mempunyai subsubbarisan
$\sequencen{x''_n}$ sedemikian rupa sehingga $\sequencen{x''_n}\halfarrow x$.

\spheader 245Yj Misalkan $\mu$ ukuran Lebesgue pada $\BbbR^r$.
Tunjukkan bahwa $L^0(\mu)$ separabel untuk topologi konvergensi dalam
ukuran. \Hint{244I.} }%akhir dari latihan

\endnotes{ \Notesheader{245} Dalam bagian ini saya mengundang Anda untuk
menganggap topologi konvergensi (lokal) dalam ukuran sebagai topologi standar
pada $L^0$, sebagaimana norma menentukan topologi standar pada ruang $L^p$
untuk $p\ge 1$. Definisi yang saya pilih dirancang agar penjumlahan dan
perkalian skalar serta operasi $\vee$, $\wedge$, dan $\times$ bersifat kontinu
(245D); lihat juga 245Xf. Dari sudut pandang analisis fungsional, sifat-sifat
ini lebih penting daripada metrizabilitas atau bahkan kelengkapan.

Seperti halnya struktur aljabar dan struktur orde pada $L^0$ dapat diuraikan
dalam teori umum ruang Riesz, hasil-hasil yang lebih mendalam 241G dan 245E juga
memiliki interpretasi dalam teori umum. Bukan kebetulan bahwa (untuk ruang
ukur semihingga) $L^0$ lengkap Dedekind jika dan hanya jika lengkap sebagai
ruang seragam; Anda dapat menemukan generalisasi yang relevan dalam 23K
dan 24E dari {\smc Fremlin 74}. Tentu saja, justru karena kedua jenis
kelengkapan ini saling terkait, saya merasa perlu menggunakan frasa
`kelengkapan Dedekind' untuk membedakan jenis kelengkapan orde khusus ini dari
kelengkapan seragam yang lebih lazim dan diuraikan dalam 2A5F.

Manfaat dari topologi konvergensi dalam ukuran sebagian besar berasal dari
245G-245J dan versi $L^p$ 245Xk dan 245Xl. Beberapa ide di sini dapat
dikaitkan dengan pertanyaan yang muncul dari teorema konvergensi dasar. Jika
$\sequencen{f_n}$ adalah suatu barisan fungsi terintegralkan yang
konvergen titik demi titik ke suatu fungsi $f$, dengan cara apa $\int f$ dapat
gagal menjadi $\lim_{n\to\infty}\int f_n$? Dalam bahasa bagian ini, hal ini
diterjemahkan menjadi: jika kita mempunyai suatu barisan (atau filter) dalam
$L^1$ yang konvergen untuk topologi konvergensi dalam ukuran, dengan cara apa
barisan ini dapat gagal konvergen dalam topologi norma pada $L^1$? Jawaban
pertama adalah Teorema Konvergensi Terdominasi Lebesgue: hal ini tidak dapat
terjadi jika barisan tersebut didominasi, yaitu, terletak dalam suatu himpunan
dari bentuk $\{u:|u|\le v\}$ di mana $v\in L^1$. (Lihat 245Xh dan 245Yg.) Saya
akan kembali ke hal ini di bagian berikutnya. Namun untuk saat ini, 245H
menunjukkan bahwa jika $\sequencen{u_n}$ konvergen dalam ukuran ke $u\in
L^1$, tetapi tidak untuk topologi norma pada $L^1$, hal itu terjadi karena
$\limsup_{n\to\infty}\|u_n\|_1$ terlalu besar; sebagian bobotnya hilang di tak
hingga, seperti dalam contoh-contoh 245I. Jika $\sequencen{u_n}$ benar-benar
ordo*-konvergen ke $u$, maka Lemma Fatou menunjukkan bahwa
$\liminf_{n\to\infty}\|u_n\|_1\ge\|u\|_1$, yaitu, bahwa fungsi limit tidak dapat
memiliki bobot lebih besar (yang diukur oleh $\|\,\|_1$) daripada bobot yang
dibawa oleh barisan. 245J dan 245Xk merupakan generalisasi fakta ini ke
konvergensi dalam ukuran. Jika Anda menginginkan generalisasi dari teorema
B. Levi, 242Yf tetap merupakan ungkapan terbaik dalam bahasa bab ini;
tetapi 245Yf adalah versi dalam istilah konsep-konsep bagian ini.

Dalam kasus ruang $\sigma$-hingga, kita mempunyai deskripsi alternatif dari
topologi konvergensi dalam ukuran (245L) yang tidak menggunakan fungsional
atau pseudometrik yang tampil dalam 245A. Hal ini dapat diungkapkan, setidaknya
dalam konteks $L^0$, dalam istilah hasil standar dari topologi umum (245Yi).
Hasil tersebut memberikan cara membentuk topologi pada
$L^0$ yang berlaku pada ruang ukur mana pun. Hal yang menarik ialah bahwa
untuk ruang $\sigma$-hingga kita mendapatkan topologi ruang linier.
}%akhir dari catatan


\discrpage
